\documentclass[12pt, a4paper]{report}

% ─── Packages ────────────────────────────────────────────────────────────────
\usepackage[top=2.5cm, bottom=2.5cm, left=3cm, right=2.5cm]{geometry}
\usepackage{graphicx}
\usepackage{xcolor}
\usepackage[table]{colortbl}
\usepackage{titlesec}
\usepackage{fancyhdr}
\usepackage{booktabs}
\usepackage{array}
\usepackage{tabularx}
\usepackage{amsmath}
\usepackage{amssymb}
\usepackage{hyperref}
\usepackage{enumitem}
\usepackage{tocloft}
\usepackage{parskip}
\usepackage{microtype}
\usepackage{setspace}
\usepackage{mdframed}
\usepackage{tikz}
\usetikzlibrary{shapes.geometric, arrows.meta, positioning, fit, backgrounds}
\usepackage{fontenc}
\usepackage{inputenc}
\usepackage{caption}

% ─── Colours ─────────────────────────────────────────────────────────────────
\definecolor{primaryblue}{RGB}{26, 60, 110}
\definecolor{accentblue}{RGB}{52, 120, 190}
\definecolor{lightgray}{RGB}{245, 247, 250}
\definecolor{darkgray}{RGB}{80, 80, 80}
\definecolor{rulecolor}{RGB}{200, 215, 230}
\definecolor{boxbg}{RGB}{235, 242, 252}

% ─── Hyperref ─────────────────────────────────────────────────────────────────
\hypersetup{
    colorlinks=true,
    linkcolor=accentblue,
    urlcolor=accentblue,
    citecolor=accentblue,
    pdftitle={Clinical Appointment No-Show Prediction},
    pdfauthor={Shaik Mohammed Junaid, Siraparapu Siva Sankar Mani Kanta, Balanagu Sesha Srikar},
}

% ─── Section Formatting ───────────────────────────────────────────────────────
\titleformat{\chapter}[display]
  {\normalfont\Large\bfseries\color{primaryblue}}
  {\chaptertitlename\ \thechapter}{12pt}{\Huge}
\titlespacing*{\chapter}{0pt}{-10pt}{20pt}

\titleformat{\section}
  {\normalfont\large\bfseries\color{primaryblue}}
  {\thesection}{1em}{}[\color{rulecolor}\titlerule]
\titlespacing*{\section}{0pt}{14pt}{6pt}

\titleformat{\subsection}
  {\normalfont\normalsize\bfseries\color{accentblue}}
  {\thesubsection}{1em}{}
\titlespacing*{\subsection}{0pt}{10pt}{4pt}

% ─── Header / Footer ─────────────────────────────────────────────────────────
\pagestyle{fancy}
\fancyhf{}
\fancyhead[L]{\small\color{darkgray}\textit{Clinical Appointment No-Show Prediction}}
\fancyhead[R]{\small\color{darkgray}\textit{Milestone 1}}
\fancyfoot[C]{\small\color{darkgray}\thepage}
\renewcommand{\headrulewidth}{0.4pt}
\renewcommand{\footrulewidth}{0pt}
\renewcommand{\headrule}{\color{rulecolor}\hrule width\headwidth height\headrulewidth}

% ─── Table of Contents Style ──────────────────────────────────────────────────
\renewcommand{\cftchapfont}{\bfseries\color{primaryblue}}
\renewcommand{\cftchappagefont}{\bfseries\color{primaryblue}}
\renewcommand{\cftsecfont}{\color{darkgray}}
\renewcommand{\cftsecpagefont}{\color{darkgray}}
\setlength{\cftbeforechapskip}{6pt}

% ─── List Styling ─────────────────────────────────────────────────────────────
\setlist[itemize]{leftmargin=1.5em, itemsep=4pt, parsep=0pt, topsep=4pt}
\setlist[enumerate]{leftmargin=1.5em, itemsep=4pt, parsep=0pt, topsep=4pt}

% ─── Line Spacing ─────────────────────────────────────────────────────────────
\setstretch{1.3}

% ─── Caption Styling ──────────────────────────────────────────────────────────
\captionsetup{font=small, labelfont={bf,color=primaryblue}, textfont=it}

% ─── Custom Boxes ─────────────────────────────────────────────────────────────
\newmdenv[
  backgroundcolor=lightgray,
  linecolor=accentblue,
  linewidth=1.5pt,
  leftline=true, rightline=false, topline=false, bottomline=false,
  innerleftmargin=10pt, innerrightmargin=10pt,
  innertopmargin=8pt, innerbottommargin=8pt,
]{infobox}

\newmdenv[
  backgroundcolor=boxbg,
  linecolor=primaryblue,
  linewidth=1pt,
  leftline=true, rightline=false, topline=false, bottomline=false,
  innerleftmargin=12pt, innerrightmargin=12pt,
  innertopmargin=10pt, innerbottommargin=10pt,
]{problembox}

% ─── TikZ styles for pipeline ─────────────────────────────────────────────────
\tikzstyle{pipenode} = [
  rectangle, rounded corners=4pt,
  minimum width=2.8cm, minimum height=0.9cm,
  text centered, text width=2.6cm,
  draw=primaryblue, fill=boxbg,
  font=\small\bfseries\color{primaryblue}
]
\tikzstyle{pipearrow} = [
  -Stealth, thick, color=accentblue
]

% ═════════════════════════════════════════════════════════════════════════════
%  DOCUMENT
% ═════════════════════════════════════════════════════════════════════════════
\begin{document}

% ─── Cover Page ──────────────────────────────────────────────────────────────
\begin{titlepage}
  \thispagestyle{empty}

  % Top accent bar
  \noindent\makebox[\linewidth]{\color{primaryblue}\rule{\paperwidth}{6pt}}
  \vspace{1.5cm}

  \begin{center}
    {\large\color{accentblue}\textsc{Introduction to Generative AI}}\\[0.4cm]
    {\normalsize\color{darkgray}Milestone 1 Report}\\[1.5cm]

    {\fontsize{28}{34}\selectfont\bfseries\color{primaryblue}
      Clinical Appointment\\[4pt]
      No-Show Prediction}\\[0.6cm]
    \color{rulecolor}\rule{0.6\linewidth}{1.5pt}\\[0.8cm]

    {\large\color{darkgray}\textit{Machine Learning Based Healthcare Operations System}}\\[2.5cm]

    % Author block — bold and colored
    \begin{tabular}{c}
      {\large\bfseries\color{primaryblue} Shaik Mohammed Junaid}\\[10pt]
      {\large\bfseries\color{primaryblue} Siraparapu Siva Sankar Mani Kanta}\\[10pt]
      {\large\bfseries\color{primaryblue} Balanagu Sesha Srikar}
    \end{tabular}\\[3cm]

    {\color{darkgray} February 2026}
  \end{center}

  \vfill
  \noindent\makebox[\linewidth]{\color{primaryblue}\rule{\paperwidth}{4pt}}
\end{titlepage}

% ─── Declaration ─────────────────────────────────────────────────────────────
\cleardoublepage
\thispagestyle{plain}
\chapter*{Declaration of Technical Integrity}
\addcontentsline{toc}{chapter}{Declaration of Technical Integrity}

\begin{infobox}
We formally affirm that the complete system architecture, problem formulation,
feature engineering strategy, model selection rationale, and evaluation
methodology were independently planned and designed by the team.

\medskip
Generative AI tools were used only in a limited and supportive capacity for
minor coding assistance, syntax refinement, and documentation formatting.
The core machine learning logic, experimental decisions, analytical reasoning,
and performance evaluation were entirely developed and validated by the team
members.
\end{infobox}

% ─── Abstract ────────────────────────────────────────────────────────────────
\cleardoublepage
\thispagestyle{plain}
\chapter*{Abstract}
\addcontentsline{toc}{chapter}{Abstract}

Missed clinical appointments significantly impact healthcare operational
efficiency, resource allocation, and patient care continuity. This project
presents a supervised machine learning framework for predicting appointment
no-shows using a dataset of 50{,}000 hospital appointment records.

Four classification algorithms---Logistic Regression, Decision Tree, Random
Forest, and XGBoost---were implemented and evaluated using Accuracy, Precision,
Recall, F1-Score, and ROC-AUC\@. Extensive feature engineering was performed
to capture behavioural history, temporal patterns, travel burden, and
communication signals.

Among the evaluated models, XGBoost achieved the highest ROC-AUC,
demonstrating strong discriminative performance. However, trade-offs between
interpretability and predictive power were carefully analysed.

% ─── Table of Contents ───────────────────────────────────────────────────────
\cleardoublepage
\tableofcontents

% ═════════════════════════════════════════════════════════════════════════════
\chapter{Introduction}
% ═════════════════════════════════════════════════════════════════════════════

Healthcare appointment no-shows lead to financial loss, inefficient scheduling,
and delayed treatment for other patients. Predicting high-risk patients enables
proactive operational interventions such as reminder systems or rescheduling
strategies.

\vspace{0.6cm}
\begin{problembox}
  \textbf{Problem Statement:} Design and implement a supervised machine
  learning system that predicts the probability of patient appointment no-shows
  using historical scheduling data.
\end{problembox}
\vspace{0.4cm}

\noindent This milestone focuses strictly on:
\begin{itemize}
  \item Data preprocessing
  \item Feature engineering
  \item Supervised classification
  \item Performance evaluation
\end{itemize}

\noindent No Large Language Models or agentic reasoning mechanisms were used in
this phase.

% ─────────────────────────────────────────────────────────────────────────────
\section{Dataset Description}
% ─────────────────────────────────────────────────────────────────────────────

This study utilises a publicly available dataset from Kaggle:
\href{https://www.kaggle.com/datasets/miadul/hospital-appointment-no-show-prediction-dataset}%
{Hospital Appointment No-Show Prediction Dataset}.
The dataset consists of \textbf{50{,}000} hospital appointment records.

\subsection*{Target Variable}
\begin{itemize}
  \item \textbf{0} --- Patient attended
  \item \textbf{1} --- Patient missed appointment
\end{itemize}

\subsection*{Feature Categories}
\begin{itemize}
  \item \textbf{Demographic:} Age, Gender, Education
  \item \textbf{Behavioural History:} Previous appointments, Previous no-shows
  \item \textbf{Temporal:} Appointment day, Lead time
  \item \textbf{Travel:} Distance, Travel time
  \item \textbf{Communication:} SMS and Email reminders
  \item \textbf{Environmental:} Rainy day, Public holiday
\end{itemize}

% ─────────────────────────────────────────────────────────────────────────────
\section{Dataset Limitations}
% ─────────────────────────────────────────────────────────────────────────────

The dataset exhibits generally weak correlations between the input features and
the target variable, indicating limited predictive strength among the available
attributes. As a result:

\begin{itemize}
  \item Most independent variables show weak associations with appointment
        no-show behaviour.
  \item Strong behavioural or temporal patterns are not clearly observable.
  \item The dataset lacks complex real-world variability such as seasonal
        effects, operational constraints, and long-term patient behavioural
        history.
  \item Feature interactions appear limited, reducing the models' ability to
        capture deeper non-linear patterns.
  \item Important factors such as socio-economic background, appointment
        urgency, prior attendance consistency, and hospital workload are not
        present.
\end{itemize}

% ═════════════════════════════════════════════════════════════════════════════
\chapter{Methodology}
% ═════════════════════════════════════════════════════════════════════════════

% ─────────────────────────────────────────────────────────────────────────────
\section{Data Preprocessing}
% ─────────────────────────────────────────────────────────────────────────────

\begin{itemize}
  \item Removal of identifier columns
  \item Handling missing values (median for numerical, mode for categorical)
  \item One-Hot Encoding for nominal categorical variables
  \item Ordinal Encoding for ordinal features with natural ordering
        (\emph{e.g.}, low $<$ medium $<$ high)
  \item Stratified 80:20 train-test split
  \item Standard scaling (Logistic Regression only)
\end{itemize}

% ─────────────────────────────────────────────────────────────────────────────
\section{Feature Engineering}
% ─────────────────────────────────────────────────────────────────────────────

\begin{itemize}
  \item No-show rate $=$ Previous no-shows $\div$ Previous appointments
  \item Lead-time indicators (short, long, same-day)
  \item Travel burden (distance $\times$ travel time)
  \item High-risk patient flag
  \item Chronic condition indicators
  \item Reminder intensity indicators
  \item Interaction features (\emph{e.g.}, uninsured $\times$ long distance)
\end{itemize}

% ─────────────────────────────────────────────────────────────────────────────
\section{Model Development}
% ─────────────────────────────────────────────────────────────────────────────

Four supervised classification algorithms were implemented:

\begin{enumerate}
  \item \textbf{Logistic Regression} --- with class-weight balancing
  \item \textbf{Decision Tree} --- with depth and leaf constraints
  \item \textbf{Random Forest} --- 200 estimators
  \item \textbf{XGBoost} --- gradient boosting with hyperparameter tuning
\end{enumerate}

% ═════════════════════════════════════════════════════════════════════════════
\chapter{Model Evaluation}
% ═════════════════════════════════════════════════════════════════════════════

% ─────────────────────────────────────────────────────────────────────────────
\section{Performance Comparison}
% ─────────────────────────────────────────────────────────────────────────────

Table~\ref{tab:performance} summarises the performance of all four models on
the held-out test set across five standard classification metrics.

\vspace{0.5cm}
\begin{table}[ht]
  \centering
  \caption{Model Performance Comparison on Hold-Out Test Set}
  \label{tab:performance}
  \renewcommand{\arraystretch}{1.8}
  \setlength{\tabcolsep}{10pt}
  \begin{tabularx}{\textwidth}{l *{5}{>{\centering\arraybackslash}X}}
    \toprule
    \rowcolor{boxbg}
    \textbf{Model} & \textbf{Accuracy} & \textbf{Precision} &
    \textbf{Recall} & \textbf{F1-Score} & \textbf{ROC-AUC} \\
    \midrule
    Logistic Regression      & 67.45\% & 89.04\% & 66.01\% & 75.82\% & 76.38\% \\
    Decision Tree            & 65.10\% & 87.28\% & 64.20\% & 73.98\% & 71.50\% \\
    Random Forest            & 71.80\% & 86.84\% & 74.86\% & 80.41\% & 75.51\% \\
    \rowcolor{boxbg}
    \textbf{XGBoost (Tuned)} & \textbf{78.43\%} & 80.18\% &
    \textbf{95.77\%} & \textbf{87.28\%} & \textbf{76.23\%} \\
    \bottomrule
  \end{tabularx}
\end{table}
\vspace{0.3cm}

\noindent The highlighted row indicates the best-performing configuration.
XGBoost achieves the highest accuracy, recall, and F1-Score, while Logistic
Regression leads on precision owing to its class-weight balancing strategy.

% ═════════════════════════════════════════════════════════════════════════════
\chapter{Trade-Off Analysis}
% ═════════════════════════════════════════════════════════════════════════════

% ─────────────────────────────────────────────────────────────────────────────
\section{Interpretability vs.\ Performance}
% ─────────────────────────────────────────────────────────────────────────────

Each model occupies a different position on the interpretability--performance
spectrum. The table below summarises the trade-off qualitatively:

\vspace{0.4cm}
\renewcommand{\arraystretch}{1.8}
\begin{tabularx}{\textwidth}{>{\bfseries\color{primaryblue}}l X X}
  \toprule
  \rowcolor{boxbg}
  \textbf{Model} & \textbf{Interpretability} & \textbf{Performance Capacity} \\
  \midrule
  Logistic Regression & High --- linear coefficients are directly readable &
    Limited non-linear modelling capacity \\
  Decision Tree       & Medium --- rule-based, human-readable paths &
    Susceptible to overfitting without constraints \\
  Random Forest       & Low --- aggregation of many trees reduces transparency &
    Improved generalisation over single trees \\
  \rowcolor{boxbg}
  XGBoost             & Lowest --- complex boosted ensemble &
    Highest predictive performance overall \\
  \bottomrule
\end{tabularx}
\vspace{0.4cm}

% ─────────────────────────────────────────────────────────────────────────────
\section{Recall vs.\ Precision}
% ─────────────────────────────────────────────────────────────────────────────

Healthcare operations prioritise high recall to minimise missed at-risk
patients. However, excessive false positives increase unnecessary interventions.
The F1-Score balances these competing objectives by computing the harmonic mean
of Precision and Recall:

\begin{equation}
  F_1 = \frac{2 \cdot \mathrm{Precision} \cdot \mathrm{Recall}}
             {\mathrm{Precision} + \mathrm{Recall}}
\end{equation}

\noindent XGBoost attains an F1-Score of \textbf{87.28\%} with a recall of
\textbf{95.77\%}, making it the most operationally suitable model for flagging
high-risk patients. Logistic Regression, despite its higher precision
(89.04\%), achieves a lower recall (66.01\%), meaning it misses a greater
proportion of actual no-show patients. The choice of operating threshold can
further tune this balance depending on clinical deployment requirements.

% ═════════════════════════════════════════════════════════════════════════════
\chapter{Feature Importance Analysis}
% ═════════════════════════════════════════════════════════════════════════════

Using XGBoost's built-in feature importance scores, the most influential
predictors for appointment no-show were identified as:

\begin{enumerate}
  \item Previous no-shows
  \item Lead time
  \item Travel distance
  \item Number of reminders
  \item Chronic condition indicators
\end{enumerate}

These findings align with domain knowledge: patients with a prior history of
missed appointments and longer waiting periods are at considerably higher risk.

% ═════════════════════════════════════════════════════════════════════════════
\chapter{System Architecture}
% ═════════════════════════════════════════════════════════════════════════════

The end-to-end pipeline was structured as a sequence of modular layers,
each with a clearly defined responsibility. Figure~\ref{fig:pipeline}
illustrates the full flow from raw data ingestion to final model comparison.

\vspace{0.6cm}

% ── TikZ Pipeline Diagram ────────────────────────────────────────────────────
\begin{figure}[ht]
\centering
\begin{tikzpicture}[node distance=0.5cm and 0.7cm]

  % Row 1
  \node (n1) [pipenode] {Dataset\\Loading};
  \node (n2) [pipenode, right=of n1] {Feature\\Engineering};
  \node (n3) [pipenode, right=of n2] {Train/Test\\Split};
  \node (n4) [pipenode, right=of n3] {Scaling \&\\Preprocessing};

  % Row 2
  \node (n5) [pipenode, below=0.9cm of n1] {Model\\Training};
  \node (n6) [pipenode, right=of n5] {Model\\Evaluation};
  \node (n7) [pipenode, right=of n6] {Performance\\Comparison};
  \node (n8) [pipenode, right=of n7] {Interpretation};

  % Row 1 arrows
  \draw [pipearrow] (n1) -- (n2);
  \draw [pipearrow] (n2) -- (n3);
  \draw [pipearrow] (n3) -- (n4);

  % Bend down from row 1 to row 2
  \draw [pipearrow] (n4.south) -- ++(0,-0.45) -| (n5.north);

  % Row 2 arrows
  \draw [pipearrow] (n5) -- (n6);
  \draw [pipearrow] (n6) -- (n7);
  \draw [pipearrow] (n7) -- (n8);

  % Background group labels
  \begin{scope}[on background layer]
    \node[draw=rulecolor, fill=white, dashed, rounded corners=6pt,
          fit=(n1)(n4), inner sep=7pt,
          label={[font=\footnotesize\color{darkgray}]above:Data Preparation Phase}] {};
    \node[draw=rulecolor, fill=white, dashed, rounded corners=6pt,
          fit=(n5)(n8), inner sep=7pt,
          label={[font=\footnotesize\color{darkgray}]below:Modelling \& Evaluation Phase}] {};
  \end{scope}

\end{tikzpicture}
\caption{End-to-end pipeline flow for the no-show prediction system}
\label{fig:pipeline}
\end{figure}

\vspace{0.4cm}

The architecture is composed of eight distinct layers:

\begin{enumerate}
  \item \textbf{Data Loading Layer} --- Ingests raw CSV records.
  \item \textbf{Feature Engineering Layer} --- Constructs derived and interaction
        features.
  \item \textbf{Data Splitting Layer} --- Performs stratified 80:20 partitioning.
  \item \textbf{Scaling Layer} --- Applies standard normalisation where required.
  \item \textbf{Preprocessing Layer} --- Encodes categorical variables.
  \item \textbf{Model Training Layer} --- Fits each classifier on the training
        partition.
  \item \textbf{Evaluation Layer} --- Computes all performance metrics.
  \item \textbf{Comparison \& Interpretation Layer} --- Benchmarks models and
        extracts feature importances.
\end{enumerate}

% ═════════════════════════════════════════════════════════════════════════════
\chapter{Results and Conclusion}
% ═════════════════════════════════════════════════════════════════════════════

% ─────────────────────────────────────────────────────────────────────────────
\section{Results}
% ─────────────────────────────────────────────────────────────────────────────

XGBoost achieved the highest ROC-AUC score, followed by Random Forest.
Logistic Regression showed competitive recall owing to class-weight balancing,
while Decision Tree provided interpretability at the cost of lower overall
performance. The full quantitative comparison is presented in
Table~\ref{tab:performance}.

% ─────────────────────────────────────────────────────────────────────────────
\section{Conclusion}
% ─────────────────────────────────────────────────────────────────────────────

This project demonstrates the effectiveness of supervised machine learning
techniques for predicting healthcare appointment no-shows. The key findings are:

\begin{itemize}
  \item \textbf{Ensemble methods outperform linear models} --- both Random Forest
        and XGBoost surpass Logistic Regression on the primary metrics.
  \item \textbf{Behavioural history is the strongest predictive signal} ---
        prior no-show rate dominates feature importance rankings.
  \item \textbf{Threshold tuning improves operational utility} --- adjusting the
        classification threshold enables a better recall--precision balance for
        real-world deployment.
  \item \textbf{Dataset limitations restrict real-world generalisability} ---
        the absence of socio-economic, urgency, and operational variables
        constrains model performance.
\end{itemize}

\noindent XGBoost provided the best overall performance and was selected as the
final model for deployment consideration in downstream milestones.

\end{document}
